\subsection{分次代数的正向极限}

设 \((\mathbf{A}_{\alpha},\phi_{\beta \alpha})\) 是 \(\Delta\) 型分次交换环的一个有向正向系统（II, § 11, 第 3 号, 注 3），且对每个 \(\alpha\) ，设 \(\mathbf{E}_{\alpha}\) 是一个 \(\mathbf{A}_{\alpha}\) 上的 \(\Delta\) 型分次代数；对 \(\alpha \leqslant \beta\) ，设 \(f_{\beta \alpha}:\mathrm{E}_{\alpha}\to \mathrm{E}_{\beta}\) 是一个 \(\mathbf{A}_{\alpha}\) -分次代数同态，并假设 \(f_{\gamma \alpha} = f_{\gamma \beta}\circ f_{\beta \alpha}\) 对 \(\alpha \leqslant \beta \leqslant \gamma\) 成立；则我们称 \((\mathrm{E}_{\alpha},f_{\beta \alpha})\) 为 \(\Delta\) 型分次交换环的有向正向系统 \((\mathbf{A}_{\alpha},\phi_{\beta \alpha})\) 上的一个 \(\Delta\) 型分次代数的有向正向系统。这时我们知道（II, § 11, 第 3 号） \(\mathbf{E} = \varinjlim \mathbf{E}_{\alpha}\) 典范地具有一个 \(\Delta\) 型分次模结构（在分次环 \(\overrightarrow{\mathbf{A}} = \varinjlim \mathbf{A}_{\alpha}\) 上），以及一个满足 \(\mathrm{E}^{\lambda}\mathrm{E}^{\mu}\subset \mathrm{E}^{\lambda +\mu}\) 的乘法（其中 \((\mathrm{E}^{\lambda})\) 表示 E 上的分次）；于是这个乘法与 E 上的分次 A-模结构在 E 上定义了一个 \(\Delta\) 型分次 A-代数结构；带有这个结构的 E 称为分次代数的正向系统 \((\mathrm{E}_{\alpha},f_{\beta \alpha})\) 的正向极限。典范同态 \(\mathrm{E}_{\alpha}\rightarrow \mathrm{E}\) 都取为 \(\mathbf{A}_{\alpha}\) -分次代数同态。此外，若 F 是 \(\Delta\) 型分次 A-代数，且 \((u_{\alpha})\) 是 \(\mathbf{A}_{\alpha}\) -同态 \(u_{\alpha}:\mathrm{E}_{\alpha}\rightarrow \mathrm{F}, u = \varinjlim u_{\alpha}\) 的一个正向系统，则 u 是分次代数的 A-同态。

